\documentclass[10pt,a4paper]{amsart}
\usepackage[margin=19mm]{geometry}
\usepackage{amsmath,amssymb,mathtools}
\usepackage[scr=rsfs,cal=euler]{mathalfa}
\usepackage{xfrac}
\usepackage[unicode,hidelinks,bookmarks=false]{hyperref}
\newcommand{\ie}{i.e.\ }
\newcommand{\cf}{cf.\ }
\newcommand{\resp}{respectively\ }
\newcommand{\Subsection}{\subsection}
\providecommand{\nobreakdash}{\nobreak\hbox{-}}
\setlength{\emergencystretch}{2em}
\multlinegap0pt
\usepackage{algorithm}\usepackage{algpseudocode}
\usepackage{stackengine}
\usepackage{amssymb}
\usepackage{mathtools}
\usepackage{xcolor}
\usepackage{siunitx}
\usepackage{bm}
\usepackage{tikz}
\newtheorem{lemma}{Lemma}
\usepackage{cleveref}
\crefname{lemma}{Lemma}{Lemmas}
\DeclareMathOperator{\Mon}{Mon}
\title{Certifying Galois/monodromy Actions via Homotopy Graphs}
\author{Timothy Duff, Kisun Lee}
\date{Source: arXiv:2603.17288v1 (CC BY 4.0)}
\begin{document}
\maketitle

\noindent\emph{Source and licence.} Certifying Galois/monodromy Actions via Homotopy Graphs. Version arXiv:2603.17288v1 (CC BY 4.0), distributed by arXiv under the Creative Commons Attribution 4.0 International licence, http://creativecommons.org/licenses/by/4.0/, as recorded in the arXiv licence metadata for this exact version and shown on the abstract page https://arxiv.org/abs/2603.17288v1. Full licence notice: \texttt{licenses/arxiv-2603.17288-CC-BY-4.0.txt}.\par\smallskip

\noindent\textit{Editorial scope.} Counted unit: the article's lemma lemma (source0.tex lines 445--447). The article's complete original proof of it is delivered with it (source0.tex lines 449--472, one \texttt{proof} environment of 24 lines). No mathematics is written, repaired or re-derived: the statement and the proof are the article's own lines, in the article's own notation, and no retained line is substituted or re-flowed.  Retained uncounted as the source-local context the proof needs: lines 409--411; lines 420--440.  Nothing was omitted from any retained span, so the package carries no omission record.  No retained line contains a citation, so the wrapper carries no bibliography and no bibliography entry.  No retained line contains a figure environment, an image-inclusion command or drawing code, so no image is delivered and the package declares no asset dependency.  Editorial formatting only, in addition: an AMS \texttt{amsart} preamble that restates the article's own declarations and macros used by the retained lines, namely the environments \texttt{lemma} on separate counters, together with the standard packages those lines need.  The retained environments print in this document as Lemma 1 in order of appearance, whereas the article prints the same items with the numbering of its own full document.  The complete native source of the article as distributed in the arXiv e-print for this exact version is bundled unchanged as \texttt{sources/196611/source0.tex}.
\begin{equation}\label{eq:tree-permutations}
\sigma_{\mathcal{T}, e} = \phi_{\mathcal{T},v}^{-1} \circ \gamma_e \circ \phi_{\mathcal{T},u},
\end{equation}

\begin{algorithm}[ht]
	\caption{MonodromyGroup}
 \label{algo:group_generation}
\begin{algorithmic}[1]
\Require  
\begin{itemize}
    \item A saturated homotopy graph $\mathcal{G} = (V, E)$ with $k$ solutions and a base vertex $z_0 \in V$,
    \item for each edge $e=(u, v) \in E$, a bijective correspondence $\gamma_e: \pi^{-1}(u) \to \pi^{-1}(v)$.
\end{itemize}
\Ensure{Generators $S \subset S_k$ for the action of $\Mon_{\mathcal{G}, z_0} \subset \Mon_\pi $ on $\pi^{-1} (z_0)$}
\State Initialize the set of generators $S = \emptyset$.
\State Construct a spanning tree $\mathcal{T}=(V,E_\mathcal{T})$ of $\mathcal{G}$ rooted at $z_0$.
\State Identify the set of off-tree edges $E_{\text{off}} = E \, \setminus \, E_\mathcal{T}$.
\For{each edge $e = (u, v) \in E_{\text{off}}$}
    \State Find $\phi_{\mathcal{T},u}$ and $\phi_{\mathcal{T},v}$ along the unique $uv$-path in $\mathcal{T}$.
    \State Construct the permutation $\sigma_{\mathcal{T},e} = \phi_{\mathcal{T},v}^{-1} \circ \gamma_e \circ \phi_{\mathcal{T},u}$ on $\pi^{-1}(z_0)$ as in~\eqref{eq:tree-permutations}.
    \State $S= S \cup \{ \sigma_{\mathcal{T},e} \}$.
\EndFor
\State \Return $S$.
 \end{algorithmic}
 \end{algorithm}

\begin{lemma}
The permutations produced by~\Cref{algo:group_generation} generate the group $\Mon_ {\mathcal{G},z_0}$.
\end{lemma}

\begin{proof}
Consider any closed walk based at $z_0,$ traversing edges $e_1=(u_1, v_1), \dots , e_k=(u_k, v_k)$, where $u_{i+1}=v_i$ for $1\le i <k,$ and $u_1 = v_k = z_0.$
This walk encodes the permutation
\begin{equation}\label{eq:perm-decomp}
\gamma_{e_k} \circ \cdots \circ \gamma_{e_1} 
=
\left( 
\phi_{\mathcal{T},v_k}^{-1}\circ  \gamma_{e_k} \circ \phi_{\mathcal{T},u_k}
\right)
\circ \cdots 
\circ
\left( 
\phi_{\mathcal{T},v_1}^{-1}\circ  \gamma_{e_1}\circ  \phi_{\mathcal{T},u_1}
\right)
.
\end{equation}
For each $e_i$, there are two cases to consider:
\begin{enumerate}
    \item $e_i \notin E_\mathcal{T}$, in which case  $\phi_{\mathcal{T},v_i}^{-1}\circ  \gamma_{e_i} \circ \phi_{\mathcal{T},u_i}= \sigma_{\mathcal{T},e_i}$ as defined in~\Cref{algo:group_generation}, or 
    
    \item $e_i\in E_\mathcal{T},$ in which case $\phi_{\mathcal{T},v_i}^{-1}\circ  \gamma_{e_i} \circ \phi_{\mathcal{T},u_i}=\text{id}.$
\end{enumerate}
We conclude that the permutation~\eqref{eq:perm-decomp} lies in the group generated by all $\sigma_{\mathcal{T},e}$ with $e\notin E_\mathcal{T}.$
\end{proof}

\end{document}
